% SPDX-License-Identifier: PMPL-1.0-or-later
% Copyright (c) 2026 Jonathan D.A. Jewell (hyperpolymath) <j.d.a.jewell@open.ac.uk>
%
% Dyadic Language Design: Co-Swimming Disciplines in Programming Language Architecture
% Target venue: arXiv (Programming Languages / PL)

\documentclass[11pt,a4paper]{article}

% --- Packages ---
\usepackage[utf8]{inputenc}
\usepackage[T1]{fontenc}
\usepackage{lmodern}
\usepackage{amsmath,amssymb,amsthm}
\usepackage{mathpartir}
\usepackage{stmaryrd}
\usepackage{xcolor}
\usepackage{hyperref}
\usepackage{cleveref}
\usepackage{booktabs}
\usepackage{enumitem}
\usepackage{listings}
\usepackage{microtype}
\usepackage[margin=2.5cm]{geometry}
\usepackage{natbib}
\usepackage{tikz}
\usetikzlibrary{arrows.meta,positioning,decorations.pathmorphing}

% --- Theorem Environments ---
\newtheorem{theorem}{Theorem}[section]
\newtheorem{lemma}[theorem]{Lemma}
\newtheorem{corollary}[theorem]{Corollary}
\newtheorem{proposition}[theorem]{Proposition}
\newtheorem{definition}[theorem]{Definition}
\newtheorem{example}[theorem]{Example}
\newtheorem{remark}[theorem]{Remark}
\newtheorem{property}[theorem]{Property}

% --- Listings Setup ---
\lstdefinelanguage{Ephapax}{
  keywords={let, let!, fn, if, then, else, match, with, region, in, drop, borrow, move, type, end, return, spawn, await},
  keywordstyle=\bfseries\color{blue!70!black},
  comment=[l]{//},
  commentstyle=\itshape\color{gray},
  stringstyle=\color{red!70!black},
  morestring=[b]",
  sensitive=true,
}
\lstset{
  basicstyle=\ttfamily\small,
  columns=flexible,
  breaklines=true,
  frame=single,
  framerule=0.4pt,
  rulecolor=\color{gray!50},
  xleftmargin=1em,
  numbers=left,
  numberstyle=\tiny\color{gray},
  numbersep=6pt,
  tabsize=2,
  captionpos=b,
}

% --- Custom Commands ---
\newcommand{\Ephapax}{\textsc{Ephapax}}
\newcommand{\letlin}{\mathbf{let!}}
\newcommand{\letaff}{\mathbf{let}}
\newcommand{\lmark}{\bullet}  % linear marker
\newcommand{\amark}{\circ}    % affine marker
\newcommand{\ctx}{\Gamma}
\newcommand{\lctx}{\Delta}    % linear sub-context
\newcommand{\judge}[3]{#1 \vdash #2 : #3}
\newcommand{\consumed}[1]{\overline{#1}}

\title{\textbf{Dyadic Language Design:\\Co-Swimming Disciplines in Programming Language Architecture}}

\author{Jonathan D.A.\ Jewell\\
  \textit{Independent Researcher}\\
  \texttt{j.d.a.jewell@open.ac.uk}}

\date{March 2026}

\begin{document}

\maketitle

% ============================================================================
% ABSTRACT
% ============================================================================
\begin{abstract}
Programming languages typically commit to a single resource-management discipline:
Rust enforces affine ownership globally, Linear Haskell annotates function arrows,
and conventional languages ignore linearity altogether. We argue that this
monoculture is a false economy. We introduce the \emph{dyadic design principle},
a methodology for constructing languages around two fundamentally different but
inseparable disciplines that co-exist at the binding level---neither diminishing
the other, together producing emergent capabilities that neither achieves alone.
We formalize three constitutive properties of dyadic pairs---\emph{non-diminishment},
\emph{emergence}, and \emph{inseparability}---and instantiate the principle in
\Ephapax{}, a language where linear (\lstinline{let!}) and affine (\lstinline{let})
bindings cohabit per-expression, enabling region-based memory management without
garbage collection and compile-time prevention of resource leaks. A case study of
56 SIGABRT crashes in a pure-affine D-Bus session manager demonstrates that the
dyadic approach prevents an entire class of failures that monodisciplinary designs
admit. We sketch three further dyadic pairs---eager/lazy evaluation,
synchronous/asynchronous control flow, and deterministic/probabilistic
computation---arguing that the principle generalizes beyond substructural types
to a broad methodology for programming language architecture.
\end{abstract}

\vspace{1em}
\noindent\textbf{Keywords:}
programming language design, substructural types, linear types, affine types,
resource management, type theory, language architecture

% ============================================================================
% 1. INTRODUCTION
% ============================================================================
\section{Introduction}
\label{sec:introduction}

The history of programming language design is, in large part, a history of
choosing sides. A language commits to garbage collection or manual memory
management; to strict evaluation or lazy evaluation; to synchronous or
asynchronous control flow. Even within the substructural type theory tradition,
languages tend toward a single dominant discipline: Rust's borrow checker
enforces an essentially affine regime across all bindings~\citep{jung2018rustbelt},
Linear Haskell restricts linearity to function arrows rather than
bindings~\citep{bernardy2018linear}, and Clean's uniqueness types impose a
uniform protocol on all unique values~\citep{barendsen1996uniqueness}.

These are not merely implementation choices---they are \emph{architectural
commitments} that shape what a language can express naturally and what it must
express awkwardly. A globally affine language like Rust requires escape hatches
(\lstinline{unsafe}, \lstinline{Rc}, \lstinline{Arc}) whenever strict
single-ownership is too rigid. A globally linear language would forbid the
implicit drops that make everyday programming ergonomic. The designer faces what
appears to be a genuine dilemma: enforce discipline and sacrifice flexibility,
or permit flexibility and sacrifice guarantees.

We contend that this dilemma is false. It arises from the unstated assumption
that a language must adopt \emph{one} discipline and apply it uniformly. This
paper introduces the \textbf{dyadic design principle}: a methodology for
constructing a language around two fundamentally different but inseparable
disciplines that operate at the same granularity---the individual binding---and
co-exist throughout the entire language.

The metaphor we adopt is biological. A mother dolphin and her calf swim in tight
formation: neither restricts the other's movement, yet their coordinated
swimming produces hydrodynamic benefits---reduced drag, improved speed---that
neither achieves alone. They are fundamentally different organisms (adult and
juvenile, teacher and learner), yet inseparable in practice during the calf's
formative period. We formalize this metaphor into three properties:

\begin{enumerate}[label=(\roman*)]
  \item \textbf{Non-diminishment.} Neither discipline restricts or weakens the
    other. The presence of affine bindings does not reduce the guarantees of
    linear bindings, and vice versa.
  \item \textbf{Emergence.} The combination produces capabilities that neither
    discipline achieves in isolation. Region-based memory management, for
    instance, emerges from the interaction of linear region handles with affine
    interior values.
  \item \textbf{Inseparability.} The disciplines are designed to co-exist, not
    to compete. Removing either discipline does not yield a ``simpler'' language
    but a \emph{crippled} one.
\end{enumerate}

We instantiate this principle in \Ephapax{}, a programming language where every
binding is introduced as either linear (\lstinline{let!}, must be consumed
exactly once) or affine (\lstinline{let}, may be dropped implicitly). This
per-binding granularity is the key architectural decision: the programmer
declares intent at the point of introduction, and the type system enforces
the corresponding discipline without global compromise.

The contributions of this paper are:

\begin{enumerate}
  \item A formal definition of the \emph{dyadic design principle} with three
    constitutive properties (\Cref{sec:dyadic-principle}).
  \item The design and formal type system of \Ephapax{}, instantiating the
    dyadic principle with linear and affine bindings at per-binding granularity
    (\Cref{sec:ephapax,sec:formal-type-system}).
  \item A region-based memory management scheme that emerges from the dyadic
    interaction of linear region handles and affine interior allocations
    (\Cref{sec:regions}).
  \item A case study of 56 SIGABRT crashes in a pure-affine D-Bus session
    manager, demonstrating that the dyadic approach prevents an entire failure
    class (\Cref{sec:case-study}).
  \item A generalization argument: three further dyadic pairs (eager/lazy,
    sync/async, deterministic/probabilistic) sketched as future instantiations
    (\Cref{sec:generalization}).
\end{enumerate}

% ============================================================================
% 2. THE DYADIC PRINCIPLE
% ============================================================================
\section{The Dyadic Principle}
\label{sec:dyadic-principle}

\subsection{Motivation: The Monodisciplinary Trap}

Consider a language designer who wishes to prevent resource leaks. The natural
impulse is to adopt linear types: every value must be consumed exactly once.
This is sound but punishing---the programmer cannot simply let a temporary
computation fall out of scope without explicitly disposing of it. The designer
then faces pressure to add implicit drops, weakening linearity to affinity.
But once affinity is the default, the ability to \emph{require} consumption
is lost, and resource leaks return through the back door.

This oscillation---between discipline and ergonomics, between safety and
usability---is the monodisciplinary trap. It arises because the designer
seeks a \emph{single} rule that governs all bindings. The dyadic principle
escapes the trap by refusing to choose.

\subsection{Formal Definition}

\begin{definition}[Dyadic Language Design]
\label{def:dyadic}
A programming language $\mathcal{L}$ exhibits \emph{dyadic design} with respect
to disciplines $D_1$ and $D_2$ if the following properties hold:

\begin{enumerate}[label=(D\arabic*)]
  \item \label{prop:nondiminishment}
    \textbf{Non-diminishment.} For every program $P$ in $\mathcal{L}$, the
    guarantees provided by $D_1$ to bindings governed by $D_1$ are identical
    whether or not $P$ also contains bindings governed by $D_2$, and
    symmetrically for $D_2$.
    \[
      \forall P.\; \text{guarantees}(D_1, P) = \text{guarantees}(D_1, P|_{D_1})
      \quad \wedge \quad
      \text{guarantees}(D_2, P) = \text{guarantees}(D_2, P|_{D_2})
    \]

  \item \label{prop:emergence}
    \textbf{Emergence.} There exists a property $\phi$ of $\mathcal{L}$ that
    holds for programs combining $D_1$ and $D_2$ bindings but does not hold
    for programs restricted to either discipline alone:
    \[
      \exists \phi.\; \mathcal{L} \models \phi
      \;\wedge\;
      \mathcal{L}|_{D_1} \not\models \phi
      \;\wedge\;
      \mathcal{L}|_{D_2} \not\models \phi
    \]

  \item \label{prop:inseparability}
    \textbf{Inseparability.} The restriction of $\mathcal{L}$ to either
    discipline alone is strictly less expressive:
    \[
      \text{Express}(\mathcal{L}|_{D_1}) \subsetneq \text{Express}(\mathcal{L})
      \quad \wedge \quad
      \text{Express}(\mathcal{L}|_{D_2}) \subsetneq \text{Express}(\mathcal{L})
    \]
\end{enumerate}
\end{definition}

The key insight is that \ref{prop:nondiminishment} prevents the disciplines
from interfering with each other, while \ref{prop:emergence} ensures that their
combination is not merely additive but \emph{synergistic}. Property
\ref{prop:inseparability} distinguishes dyadic design from the trivial case of
two unrelated features bolted onto a language: removing either discipline must
cause genuine loss.

\subsection{The Co-Swimming Metaphor}

The dolphin metaphor is not merely illustrative---it captures the essential
topology of the design. Two dolphins in echelon formation experience reduced
hydrodynamic drag through wave interference effects~\citep{weihs2004mechanics}.
The key properties map precisely:

\begin{itemize}
  \item \textbf{Never less than one.} Each dolphin swims with its full
    capability. The calf is not ``carried'' by the mother; it swims
    independently but in coordination.
  \item \textbf{Together more than two.} The drag reduction is a property of
    the \emph{pair}, not of either individual. It emerges from their relative
    positioning.
  \item \textbf{Fundamentally different.} The adult and calf have different
    sizes, speeds, and experience levels. The coordination protocol accommodates
    this asymmetry.
  \item \textbf{Inseparable in practice.} Separating the pair does not yield
    two independent swimmers with unchanged properties; it yields a slower calf
    and a mother who must expend energy searching.
\end{itemize}

In language design terms: each discipline operates at full strength on its own
bindings (\emph{never less than one}); the combination enables idioms impossible
with either alone (\emph{together more than two}); the disciplines differ in
kind, not merely in degree (\emph{fundamentally different}); and removing either
cripples the language (\emph{inseparable}).

\subsection{Distinguishing Dyadic Design from Multi-Paradigm Languages}

One might object that multi-paradigm languages (Scala, OCaml, Common Lisp)
already combine multiple disciplines. However, multi-paradigm languages
typically violate \ref{prop:nondiminishment}: the object-oriented features of
Scala interfere with its functional features (mutable state undermines
referential transparency), and the dynamic features of Common Lisp undermine
its static type declarations. In a multi-paradigm language, the disciplines
coexist but do not \emph{co-swim}---they compete for dominance, and programs
tend to gravitate toward one paradigm or the other.

Dyadic design, by contrast, requires that both disciplines be present in
every non-trivial program and that neither corrode the other's guarantees.
The disciplines are not alternative paradigms but complementary aspects of
a single coherent design.

% ============================================================================
% 3. BACKGROUND
% ============================================================================
\section{Background}
\label{sec:background}

\subsection{Linear Logic and Substructural Type Systems}

Girard's linear logic~\citep{girard1987linear} introduced the principle that
propositions---and by the Curry-Howard correspondence, values---are resources
that must be used exactly once. The structural rules of classical logic that
linear logic rejects are:

\begin{itemize}
  \item \textbf{Weakening} (discarding unused hypotheses): $\ctx, A \vdash C$
    implies $\ctx \vdash C$.
  \item \textbf{Contraction} (duplicating hypotheses): $\ctx, A, A \vdash C$
    implies $\ctx, A \vdash C$.
\end{itemize}

A \emph{linear} type system forbids both weakening and contraction: every
linear value must be used exactly once. An \emph{affine} type system forbids
contraction but permits weakening: affine values may be used at most once (they
can be silently dropped). A \emph{relevant} type system forbids weakening but
permits contraction: relevant values must be used at least once. These form
a lattice of substructural disciplines~\citep{walker2005substructural}.

\subsection{Rust: Global Affinity}

Rust's ownership system is essentially an affine type system applied
globally~\citep{jung2018rustbelt}. Every value has a single owner, ownership
can be transferred (moved), and values are dropped when their owner goes out of
scope. This is affine: values are used at most once (moved or dropped), but the
programmer cannot express that a value \emph{must} be consumed.

Rust compensates with runtime mechanisms: \lstinline{Drop} traits provide
destructors, \lstinline{#[must_use]} attributes generate warnings (not errors),
and the \lstinline{ManuallyDrop} wrapper suppresses the implicit destructor. But
these are advisory, not structural. A \lstinline{#[must_use]} value can still be
silently bound to \lstinline{_} and dropped. The language has no way to express
``this file handle \emph{must} be closed by the programmer, not by the
destructor.''

\subsection{Linear Haskell: Arrow-Level Linearity}

Linear Haskell~\citep{bernardy2018linear} introduces linearity at the level of
function arrows: $A \multimap B$ means that if the function's result is consumed
exactly once, then its argument is consumed exactly once. This is a weaker
commitment than per-binding linearity---the linearity is contingent on the
calling context rather than declared at the binding site.

Linear Haskell's approach is elegant but limited. Because linearity is a
property of arrows rather than bindings, it cannot express that a particular
\lstinline{let}-binding within a function body must be consumed. The linearity
``flows through'' the arrow but does not attach to intermediate values.

\subsection{Other Substructural Languages}

\textbf{ATS}~\citep{xi2017applied} combines dependent types with linear types,
using ``views'' (linear propositions about memory) to achieve safe manual memory
management. ATS is arguably the closest precursor to dyadic design, but its
linear and non-linear worlds are separated by proof-level/value-level
stratification rather than unified at the binding level.

\textbf{Granule}~\citep{orchard2019quantitative} implements quantitative type
theory (QTT) where resource usage is tracked by a semiring of grades. This
subsumes both linear and affine types as special cases ($1$ for linear,
$\{0,1\}$ for affine). Granule's approach is more general than dyadic design
but also more complex: the programmer must reason about an algebra of grades
rather than a binary choice.

\textbf{Idris 2}~\citep{brady2021idris2} also uses QTT, tracking usage as
$0$ (erased), $1$ (linear), or $\omega$ (unrestricted). Like Granule, this
subsumes linear and affine types but at the cost of a more complex mental model.

\textbf{Clean}~\citep{barendsen1996uniqueness} uses uniqueness types, which are
related to linear types but focus on reference uniqueness rather than usage
counting. A unique value is guaranteed to have no other references, enabling
in-place mutation.

% ============================================================================
% 4. EPHAPAX: A DYADIC LANGUAGE
% ============================================================================
\section{Ephapax: A Dyadic Language}
\label{sec:ephapax}

\Ephapax{}\footnote{From Greek \textit{ephapax} (\greektext{\'ef'apax}\latintext{}),
``once for all''---a term evoking the finality of linear consumption.} is a
programming language designed from the ground up around the dyadic principle,
instantiated with linear and affine types as the co-swimming disciplines. The
central architectural decision is \textbf{per-binding discipline choice}: every
binding is introduced as either linear or affine, and the type system enforces
the corresponding invariant.

\subsection{Binding Forms}

\Ephapax{} provides two binding forms:

\begin{lstlisting}[language=Ephapax, caption={The two binding forms of \Ephapax{}.}]
// Affine binding: may be dropped implicitly
let x = expensive_computation()

// Linear binding: MUST be consumed exactly once
let! handle = open_file("data.csv")
\end{lstlisting}

The \lstinline{let!} form introduces a \emph{linear} binding: the variable
\lstinline{handle} must appear exactly once in the continuation, either consumed
by a function or explicitly transferred. Failure to consume a linear binding, or
consuming it more than once, is a \emph{compile-time error}.

The \lstinline{let} form introduces an \emph{affine} binding: the variable
\lstinline{x} may be used at most once and may be silently dropped when it goes
out of scope. This provides the ergonomic default that everyday programming
requires.

\subsection{Why Per-Binding Granularity Matters}

The per-binding choice is not merely syntactic sugar. It reflects a fundamental
design decision about \emph{where discipline is declared}. Consider the
alternatives:

\begin{itemize}
  \item \textbf{Per-type} (as in Rust): All values of type \lstinline{File}
    have the same discipline. This conflates the type's nature with its usage
    protocol. A \lstinline{File} that is opened for logging (write-and-forget)
    has different lifecycle needs than a \lstinline{File} opened for a database
    transaction (must-close-before-commit).

  \item \textbf{Per-arrow} (as in Linear Haskell): Discipline is declared on
    functions, not values. This makes it impossible to express that one local
    variable must be consumed while another in the same function body may be
    dropped.

  \item \textbf{Per-binding} (as in \Ephapax{}): Discipline is declared at the
    point of introduction. The programmer knows, when writing the binding, what
    the lifecycle contract is. The type system enforces it from that point forward.
\end{itemize}

Per-binding granularity is the natural home for the dyadic principle because
it allows both disciplines to coexist within a single expression, a single
function body, even a single \lstinline{match} arm. The disciplines swim
together at the finest grain.

\subsection{Explicit Borrowing}

\Ephapax{} provides explicit borrowing via the \lstinline{&} operator:

\begin{lstlisting}[language=Ephapax, caption={Borrowing provides temporary read access without consuming.}]
let! connection = open_database("postgres://...")

// Borrow for read access: does not consume the connection
let status = check_health(&connection)

// connection is still live and must be consumed
close_database(connection)
\end{lstlisting}

A borrow \lstinline{&x} creates a temporary, non-owning reference to \lstinline{x}.
Borrows are always affine (they may be dropped) regardless of whether the
underlying binding is linear or affine. This is a deliberate asymmetry: the
borrow \emph{does not consume} the original, so it must not itself carry a
consumption obligation.

Borrowing satisfies non-diminishment: borrowing a linear value does not weaken
its linearity. The linear binding must still be consumed exactly once; the
borrow is a separate, ephemeral view.

\subsection{A Complete Example}

\begin{lstlisting}[language=Ephapax, caption={Dyadic resource management in \Ephapax{}.}, label=lst:complete]
fn transfer_funds(from_id: AccountId, to_id: AccountId, amount: Money) {
  // Linear: database connection MUST be properly closed
  let! db = open_connection("postgres://bank")

  // Linear: transaction MUST be committed or rolled back
  let! txn = begin_transaction(&db)

  // Affine: intermediate query results can be dropped
  let from_balance = query_balance(&txn, from_id)
  let to_balance = query_balance(&txn, to_id)

  if from_balance >= amount then
    execute(&txn, debit(from_id, amount))
    execute(&txn, credit(to_id, amount))
    // Consume txn: commits the transaction
    commit(txn)
  else
    // Consume txn: rolls back the transaction
    rollback(txn)
  end

  // Consume db: closes the connection
  close_connection(db)
}
\end{lstlisting}

In \Cref{lst:complete}, the linear bindings \lstinline{db} and \lstinline{txn}
guarantee at compile time that every control-flow path either commits or rolls
back the transaction and closes the connection. The affine bindings
\lstinline{from_balance} and \lstinline{to_balance} may be silently dropped
because their lifecycle is unimportant---they are intermediate computations,
not resources with external obligations.

This is the dyadic principle in action: the linear discipline ensures resource
safety; the affine discipline ensures ergonomic everyday programming; and their
coexistence within a single function body is seamless.

% ============================================================================
% 5. FORMAL TYPE SYSTEM
% ============================================================================
\section{Formal Type System}
\label{sec:formal-type-system}

\subsection{Contexts and Qualifiers}

We define a typing context $\ctx$ as a sequence of bindings $x :^q A$ where
$q \in \{\lmark, \amark\}$ is a \emph{qualifier}: $\lmark$ for linear, $\amark$
for affine.

\begin{definition}[Context Well-formedness]
A context $\ctx$ is \emph{well-formed} if no variable appears more than once
and every binding carries a qualifier:
\[
  \ctx ::= \emptyset \mid \ctx, x :^q A \qquad (q \in \{\lmark, \amark\})
\]
\end{definition}

\begin{definition}[Context Split]
A context split $\ctx = \ctx_1 \bowtie \ctx_2$ partitions $\ctx$ such that:
\begin{itemize}
  \item Every linear binding $x :^\lmark A \in \ctx$ appears in exactly one of
    $\ctx_1, \ctx_2$.
  \item Every affine binding $x :^\amark A \in \ctx$ appears in at most one of
    $\ctx_1, \ctx_2$ (and may appear in neither, i.e., it may be weakened away).
\end{itemize}
\end{definition}

Context splitting is the mechanism by which the type system tracks resource
usage: when a term has two sub-terms, the context is split between them, and the
split respects the qualifier of each binding.

\subsection{Typing Rules}

The core typing rules of \Ephapax{} are given in \Cref{fig:typing-rules}.

\begin{figure}[htbp]
\centering
\begin{mathpar}
  \inferrule*[right=\textsc{Var}]
    {~}
    {\judge{x :^q A}{x}{A}}

  \inferrule*[right=\textsc{Let-Affine}]
    {\judge{\ctx_1}{e_1}{A} \\ \judge{\ctx_2, x :^\amark A}{e_2}{B}}
    {\judge{\ctx_1 \bowtie \ctx_2}{\letaff\; x = e_1 \;\mathbf{in}\; e_2}{B}}

  \inferrule*[right=\textsc{Let-Linear}]
    {\judge{\ctx_1}{e_1}{A} \\ \judge{\ctx_2, x :^\lmark A}{e_2}{B}
     \\ x \in \text{FV}(e_2)}
    {\judge{\ctx_1 \bowtie \ctx_2}{\letlin\; x = e_1 \;\mathbf{in}\; e_2}{B}}

  \inferrule*[right=\textsc{Abs}]
    {\judge{\ctx, x :^q A}{e}{B}}
    {\judge{\ctx}{\lambda x :^q A.\; e}{A \xrightarrow{q} B}}

  \inferrule*[right=\textsc{App}]
    {\judge{\ctx_1}{e_1}{A \xrightarrow{q} B} \\ \judge{\ctx_2}{e_2}{A}}
    {\judge{\ctx_1 \bowtie \ctx_2}{e_1\; e_2}{B}}

  \inferrule*[right=\textsc{Borrow}]
    {\judge{\ctx}{e}{A} \\ x :^q A \in \ctx}
    {\judge{\ctx}{\&x}{\&A}}

  \inferrule*[right=\textsc{Weaken}]
    {\judge{\ctx}{e}{B} \\ x :^\amark A \notin \text{FV}(e)}
    {\judge{\ctx, x :^\amark A}{e}{B}}

  \inferrule*[right=\textsc{Branch}]
    {\judge{\ctx}{e_0}{\text{Bool}} \\
     \judge{\ctx_1}{e_1}{A} \\
     \judge{\ctx_2}{e_2}{A} \\
     \ctx_1 \sim_\lmark \ctx_2}
    {\judge{\ctx \bowtie \ctx_1 \bowtie \ctx_2}
           {\mathbf{if}\; e_0\; \mathbf{then}\; e_1\; \mathbf{else}\; e_2}{A}}
\end{mathpar}
\caption{Core typing rules of \Ephapax{}. The notation $\ctx_1 \sim_\lmark \ctx_2$
means that $\ctx_1$ and $\ctx_2$ consume the same set of linear bindings.}
\label{fig:typing-rules}
\end{figure}

The critical rules are:

\begin{itemize}
  \item \textsc{Let-Linear} requires that $x$ appears free in $e_2$---the linear
    binding must be used. This is the ``must consume'' invariant.
  \item \textsc{Weaken} permits affine bindings to be silently dropped. This rule
    applies \emph{only} to affine-qualified bindings.
  \item \textsc{Branch} requires that both branches consume the same linear
    bindings ($\ctx_1 \sim_\lmark \ctx_2$). This prevents a linear resource
    from being consumed in one branch but leaked in another.
\end{itemize}

\subsection{Non-Diminishment in the Type System}

We can now state \ref{prop:nondiminishment} precisely:

\begin{theorem}[Non-Diminishment]
\label{thm:nondiminishment}
Let $P$ be a well-typed \Ephapax{} program and let $x :^\lmark A$ be a linear
binding in $P$. Then $x$ is consumed exactly once in $P$, regardless of the
qualifiers of all other bindings in $P$. Symmetrically, if $y :^\amark B$ is
an affine binding, then $y$ is consumed at most once, regardless of the
qualifiers of all other bindings.
\end{theorem}

\begin{proof}[Proof sketch]
By induction on the typing derivation. The context split $\bowtie$ partitions
bindings independently: the split decision for a linear binding $x$ depends only
on whether $x$ is used in the left or right sub-term, not on the qualifiers of
other bindings. The \textsc{Weaken} rule applies only to affine bindings and
does not affect linear bindings in the context. The \textsc{Branch} rule
enforces uniform linear consumption across branches independently of affine
bindings. Thus, changing the qualifier of one binding does not affect the
enforcement of another binding's discipline. \qed
\end{proof}

\subsection{Type Soundness}

We establish type soundness via the standard progress and preservation theorems.

\begin{definition}[Values and Reduction]
Values $v$ include literals, closures $\langle \lambda x :^q A.\; e, \rho \rangle$,
and region handles. The small-step reduction relation $e \longrightarrow e'$
includes beta-reduction, let-binding elimination, branch reduction, and region
operations.
\end{definition}

\begin{theorem}[Progress]
\label{thm:progress}
If $\judge{\emptyset}{e}{A}$ then either $e$ is a value or there exists $e'$
such that $e \longrightarrow e'$.
\end{theorem}

\begin{proof}[Proof sketch]
By case analysis on the typing derivation. In the empty context, all bindings
have been provided, so every eliminator (application, branch, match) has its
principal argument available. \qed
\end{proof}

\begin{theorem}[Preservation]
\label{thm:preservation}
If $\judge{\ctx}{e}{A}$ and $e \longrightarrow e'$ then $\judge{\ctx'}{e'}{A}$
where $\ctx'$ is obtained from $\ctx$ by consuming bindings used in the
reduction step, respecting their qualifiers.
\end{theorem}

\begin{proof}[Proof sketch]
By case analysis on the reduction step, using the substitution lemma: if
$\judge{\ctx, x :^q A}{e}{B}$ and $\judge{\ctx'}{v}{A}$ then
$\judge{\ctx \bowtie \ctx'}{e[v/x]}{B}$. The qualifier is preserved through
substitution because the split respects $q$. \qed
\end{proof}

\begin{corollary}[No Stuck States]
A well-typed closed \Ephapax{} program does not get stuck: it either reduces
to a value or diverges.
\end{corollary}

\subsection{The Linear Safety Theorem}

Beyond standard soundness, we prove a property specific to the dyadic design:

\begin{theorem}[Linear Safety]
\label{thm:linear-safety}
If $\judge{\emptyset}{e}{A}$ and $e \longrightarrow^* v$ (i.e., $e$ evaluates
to a value), then every linear binding $x :^\lmark B$ introduced during
evaluation is consumed exactly once.
\end{theorem}

\begin{proof}[Proof sketch]
By \Cref{thm:preservation}, each reduction step preserves the typing invariant.
The \textsc{Let-Linear} rule ensures that $x$ appears in $\text{FV}(e_2)$, so
the binding will eventually be consumed. Context splitting ensures it is not
consumed more than once. By induction on the evaluation length, every linear
binding introduced is eventually consumed exactly once. \qed
\end{proof}

This theorem is the formal counterpart of the practical guarantee: resources
bound with \lstinline{let!} cannot leak.

% ============================================================================
% 6. REGION-BASED MEMORY MANAGEMENT
% ============================================================================
\section{Region-Based Memory Management}
\label{sec:regions}

The dyadic interaction between linear and affine bindings enables a
region-based memory management scheme that requires no garbage collector and
no manual \lstinline{free} calls.

\subsection{Regions as Linear Handles}

A \emph{region} is an arena allocator represented by a linear handle:

\begin{lstlisting}[language=Ephapax, caption={Region-based allocation.}]
fn process_batch(items: List<Item>) {
  // Linear: region MUST be explicitly closed (memory freed)
  let! arena = region::create()

  // Affine: allocations within the region are individually droppable
  let buffer = region::alloc(&arena, 1024)
  let index  = region::alloc(&arena, build_index(items))

  // ... computation using buffer and index ...

  // Consume arena: all memory in the region is freed at once
  region::destroy(arena)
}
\end{lstlisting}

The region handle \lstinline{arena} is linear: it must be consumed by
\lstinline{region::destroy}, which frees all memory allocated within it.
Individual allocations (\lstinline{buffer}, \lstinline{index}) are affine: they
may be dropped before the region is destroyed, but their memory is not freed
until the region itself is consumed.

\subsection{Emergence: The Region Safety Property}

This is a concrete instance of \ref{prop:emergence}:

\begin{theorem}[Region Safety]
\label{thm:region-safety}
In a well-typed \Ephapax{} program:
\begin{enumerate}
  \item Every region that is created is eventually destroyed (by Linear Safety,
    \Cref{thm:linear-safety}).
  \item No allocation within a region is accessed after the region is destroyed
    (by the borrow checker: borrows of region-interior values cannot outlive the
    region handle).
  \item Memory is freed exactly once, when the region handle is consumed.
\end{enumerate}
\end{theorem}

This property is \emph{emergent}: it does not hold in a purely linear language
(where interior allocations would also need explicit deallocation, making arenas
pointless) or in a purely affine language (where the region handle could be
silently dropped, leaking all its memory). It arises only from the
\emph{interaction} of linear handles with affine interiors.

\subsection{Lifetime Proofs}

Borrows within a region carry lifetime annotations that are checked at compile
time:

\begin{lstlisting}[language=Ephapax, caption={Lifetime-checked borrows.}]
fn bad_example() {
  let! arena = region::create()
  let data = region::alloc(&arena, 42)

  // ERROR: borrow of data escapes the region lifetime
  let escaped = &data
  region::destroy(arena)

  // escaped would be a dangling reference here
  // Compile-time error: borrow outlives region
}
\end{lstlisting}

The lifetime system ensures that borrows of region-interior values do not
outlive the region handle. Since the handle is linear and must be consumed,
there is a well-defined point at which the region ceases to exist, and all
borrows must expire before that point.

\subsection{Comparison with Rust's Lifetimes}

Rust's lifetime system achieves a similar effect but through a different
mechanism: all values are affine (implicitly droppable), and lifetimes are
enforced by the borrow checker. \Ephapax{}'s approach is more explicit: the
linear handle announces ``I am a resource with a mandatory lifecycle,'' and the
affine interiors announce ``I am a value whose lifecycle is subordinate to the
region.'' The intent is declared, not inferred.

This explicitness has a cost (the programmer writes \lstinline{let!} for
resources) but also a benefit: the code is self-documenting about which values
are resources and which are data. A reader of \Cref{lst:complete} can
immediately see that \lstinline{db} and \lstinline{txn} are resources with
obligations, while \lstinline{from_balance} is ephemeral.

% ============================================================================
% 7. CASE STUDY: SESSION-SENTINEL
% ============================================================================
\section{Case Study: Session-Sentinel Failure Analysis}
\label{sec:case-study}

\subsection{Background}

\textit{Session-sentinel} is a D-Bus session manager for Linux desktop
environments, originally implemented in a language with pure-affine semantics
(all values implicitly droppable, no linear binding discipline). Over a period
of several weeks of continuous deployment, the system produced 56 SIGABRT
crashes---abrupt, unrecoverable process terminations triggered by the D-Bus
library when connection handles were leaked.

\subsection{Root Cause}

The SIGABRT crashes shared a common pattern: a D-Bus connection handle was
obtained, used for one or more method calls, and then silently dropped when
it went out of scope. The D-Bus library's internal state machine expected the
connection to be explicitly closed (\lstinline{dbus_connection_close}) before
the handle was freed (\lstinline{dbus_connection_unref}). When a handle was
dropped without closing, the library detected the inconsistency and called
\lstinline{abort()}.

In pseudocode:

\begin{lstlisting}[language=Ephapax, caption={The pure-affine pattern that caused SIGABRT.}]
fn handle_session_event(event: SessionEvent) {
  // Affine: can be dropped implicitly
  let conn = dbus_connect("org.freedesktop.login1")

  match event with
  | Suspend  -> dbus_call(&conn, "Suspend")
  | Shutdown -> dbus_call(&conn, "PowerOff")
  | Lock     -> dbus_call(&conn, "LockSessions")
  end

  // conn goes out of scope here
  // Affine semantics: silently dropped
  // D-Bus library: SIGABRT! Connection not properly closed!
}
\end{lstlisting}

The programmer's intent was correct---the connection should be closed after
use---but the language provided no mechanism to \emph{enforce} that intent.
The affine type system allowed the handle to be dropped, and the drop path
did not include the required \lstinline{dbus_connection_close} call.

\subsection{The Dyadic Solution}

In \Ephapax{}, the connection handle would be bound linearly:

\begin{lstlisting}[language=Ephapax, caption={The dyadic pattern that prevents SIGABRT.}]
fn handle_session_event(event: SessionEvent) {
  // Linear: MUST be consumed (closed) explicitly
  let! conn = dbus_connect("org.freedesktop.login1")

  match event with
  | Suspend  -> dbus_call(&conn, "Suspend")
  | Shutdown -> dbus_call(&conn, "PowerOff")
  | Lock     -> dbus_call(&conn, "LockSessions")
  end

  // Consume conn: properly closes the connection
  dbus_close(conn)
}
\end{lstlisting}

If the programmer omits the \lstinline{dbus_close(conn)} line, the compiler
rejects the program: ``linear binding \lstinline{conn} is not consumed.'' The
56 SIGABRT crashes are prevented at compile time, with zero runtime cost.

\subsection{Why Pure Linear Types Are Not the Answer}

One might argue that the solution is simply to make \emph{all} types linear,
avoiding the need for the dyadic approach. But this is precisely the
monodisciplinary trap. In the session-sentinel code, many values are genuinely
ephemeral: event payloads, string formatting results, intermediate computations.
Forcing all of these to be explicitly consumed would make the code verbose and
fragile:

\begin{lstlisting}[language=Ephapax, caption={The purely-linear alternative: verbose and fragile.}]
fn handle_session_event(event: SessionEvent) {
  let! conn = dbus_connect("org.freedesktop.login1")
  let! event_str = format_event(event)  // Must consume!
  let! log_result = log_info(event_str) // Must consume!
  let! _unit = consume(log_result)       // Boilerplate disposal

  // ... actual logic ...

  dbus_close(conn)
}
\end{lstlisting}

The dyadic approach avoids this: \lstinline{event_str} and
\lstinline{log_result} are affine (\lstinline{let}), so they can be dropped
silently. Only the resource that \emph{matters}---the D-Bus connection---is
linear. The programmer expresses intent, and the type system enforces it,
without imposing discipline where it is not needed.

\subsection{Quantitative Analysis}

Of the 56 SIGABRT crashes:

\begin{table}[htbp]
\centering
\begin{tabular}{lrr}
\toprule
\textbf{Category} & \textbf{Count} & \textbf{Prevented by} \lstinline{let!}? \\
\midrule
Connection handle leaked on normal path & 31 & Yes \\
Connection handle leaked on error path  & 18 & Yes \\
Connection handle leaked in timeout     & 5  & Yes \\
Connection double-free (ref counting)   & 2  & Yes (linearity prevents aliasing) \\
\bottomrule
\end{tabular}
\caption{Classification of session-sentinel SIGABRT crashes.}
\label{tab:crashes}
\end{table}

All 56 crashes would be prevented by binding the connection handle with
\lstinline{let!}. The \textsc{Branch} typing rule (\Cref{fig:typing-rules})
ensures that all control paths---including error paths and timeouts---consume
the linear handle, preventing the error-path leaks that accounted for 18 of the
56 crashes.

% ============================================================================
% 8. WEBASSEMBLY COMPILATION TARGET
% ============================================================================
\section{Compilation to WebAssembly}
\label{sec:wasm}

\Ephapax{} targets WebAssembly (Wasm) as its primary compilation backend. The
dyadic type system enables optimizations at the Wasm level that would be
unavailable in a monodisciplinary language.

\subsection{Linear Bindings and Stack Allocation}

Because linear bindings have a statically known single-use point, the compiler
can allocate them on the Wasm stack rather than the linear memory heap.
The binding is produced, consumed, and its stack slot is immediately
reclaimable. This eliminates allocation overhead for linear values entirely.

\subsection{Affine Bindings and Arena Compaction}

Affine bindings that are allocated within a region can be placed in a compact
arena in Wasm linear memory. Because they are individually droppable but
collectively freed with the region, the allocator does not need to track
individual deallocations---it simply bumps a pointer for allocation and resets
the arena pointer for deallocation.

\subsection{No Runtime Overhead for Qualifiers}

The qualifiers $\lmark$ and $\amark$ exist only at compile time. At the Wasm
level, linear and affine values are represented identically. The type system's
guarantees are enforced entirely by the compiler, with zero runtime overhead.
This is a key advantage over garbage-collected languages and over Rust's
\lstinline{Drop}-trait-based resource management, which requires runtime
destructor dispatch.

% ============================================================================
% 9. GENERALIZATION: OTHER DYADIC PAIRS
% ============================================================================
\section{Generalization: Other Dyadic Pairs}
\label{sec:generalization}

The dyadic design principle is not specific to linear and affine types. We
sketch three further instantiations, each of which satisfies the three
constitutive properties.

\subsection{Eager/Lazy Evaluation}

Most languages commit to either strict (eager) evaluation (ML, Rust, Java) or
non-strict (lazy) evaluation (Haskell, Clean). A dyadic language could provide
both at the binding level:

\begin{lstlisting}[language=Ephapax, caption={Hypothetical eager/lazy dyadic bindings.}]
// Eager: evaluated immediately, value available now
let! result = expensive_computation()

// Lazy: thunk, evaluated on first access
let~ deferred = expensive_computation()
\end{lstlisting}

\begin{itemize}
  \item \textbf{Non-diminishment:} Eager bindings evaluate with the same
    semantics regardless of the presence of lazy bindings, and vice versa.
  \item \textbf{Emergence:} The combination enables optimistic/speculative
    patterns: eager bindings for values known to be needed, lazy bindings for
    values that may not be. The interaction between eager producers and lazy
    consumers enables data-flow patterns that neither evaluation strategy
    achieves alone.
  \item \textbf{Inseparability:} A purely eager language wastes work on unused
    computations; a purely lazy language introduces unpredictable latency. The
    dyadic combination avoids both failure modes.
\end{itemize}

This dyadic pair would subsume Haskell's \lstinline{seq}/\lstinline{BangPatterns}
escape hatches and Scala's \lstinline{lazy val}, but as a first-class
architectural principle rather than a bolt-on pragmatic compromise.

\subsection{Synchronous/Asynchronous Control Flow}

The sync/async dichotomy is one of the most painful in modern language design.
Rust's colored function problem~\citep{nystrom2015colored}, JavaScript's
\lstinline{async/await} proliferation, and Go's goroutine model all represent
different monodisciplinary commitments. A dyadic approach:

\begin{lstlisting}[language=Ephapax, caption={Hypothetical sync/async dyadic bindings.}]
// Synchronous: blocks until complete
let! data = read_file("config.toml")

// Asynchronous: returns a future, non-blocking
let& task = fetch_url("https://api.example.com/data")

// Synchronous consumption of async result
let! result = await(task)
\end{lstlisting}

\begin{itemize}
  \item \textbf{Non-diminishment:} Synchronous bindings block as expected,
    unaffected by the presence of asynchronous bindings in the same scope.
    Asynchronous bindings are non-blocking regardless of surrounding synchronous
    code.
  \item \textbf{Emergence:} The type system can enforce that synchronous code
    does not accidentally depend on an unawaited future (a common bug in
    JavaScript and Python). The combination enables structured concurrency
    where the boundaries between sync and async are explicit and
    type-checked.
  \item \textbf{Inseparability:} Purely synchronous code cannot express
    concurrent operations; purely asynchronous code (as in early Node.js)
    forces callback indirection even for naturally sequential operations.
\end{itemize}

The key insight is that the ``function color'' problem arises because
sync/async is treated as a \emph{type-level} distinction (colored functions)
rather than a \emph{binding-level} one. Moving the distinction to the binding
level---where the dyadic principle naturally operates---dissolves the problem.

\subsection{Deterministic/Probabilistic Computation}

A more speculative instantiation: a language for probabilistic programming
where deterministic and probabilistic bindings co-swim:

\begin{lstlisting}[language=Ephapax, caption={Hypothetical deterministic/probabilistic dyadic bindings.}]
// Deterministic: exact value, no uncertainty
let! pi = 3.14159265358979

// Probabilistic: distribution over values
let~ temperature = normal(20.0, 2.0)

// Deterministic function of a probabilistic value:
// result is probabilistic (discipline propagates)
let~ adjusted = temperature + 5.0

// Deterministic observation: collapses probability
let! sample = observe(temperature, 21.5)
\end{lstlisting}

\begin{itemize}
  \item \textbf{Non-diminishment:} Deterministic computations are exact
    regardless of probabilistic bindings in scope. Probabilistic bindings
    maintain their distributional semantics regardless of deterministic
    context.
  \item \textbf{Emergence:} Bayesian inference emerges naturally from the
    interaction: \lstinline{observe} collapses a probabilistic binding into
    a deterministic one, updating the posterior distribution of related
    probabilistic bindings. This is a property of the \emph{pair}, not of
    either discipline alone.
  \item \textbf{Inseparability:} A purely deterministic language cannot
    express uncertainty; a purely probabilistic language makes exact computation
    awkward (everything becomes a point distribution).
\end{itemize}

\subsection{Toward a Taxonomy of Dyadic Pairs}

The four dyadic pairs we have discussed---linear/affine, eager/lazy,
sync/async, deterministic/probabilistic---suggest a pattern. Each pair consists
of:

\begin{enumerate}
  \item A \textbf{strict discipline} that imposes obligations (linear, eager,
    synchronous, deterministic).
  \item A \textbf{permissive discipline} that relaxes obligations (affine, lazy,
    asynchronous, probabilistic).
\end{enumerate}

This is not a coincidence. The dyadic principle is, at its core, about the
productive tension between obligation and freedom. A language that provides
only obligation is unusable; a language that provides only freedom is unsafe.
The dyadic principle says: provide both, at the binding level, and let the
programmer declare which regime each value lives under.

We conjecture that any productive tension in language design that manifests as
a global architectural choice (``strict or lazy?'', ``sync or async?'') can be
refactored into a dyadic pair at the binding level. Testing this conjecture
across further domains is future work.

% ============================================================================
% 10. RELATED WORK
% ============================================================================
\section{Related Work}
\label{sec:related-work}

\paragraph{Rust and Ownership.}
Rust~\citep{jung2018rustbelt} is the most widely deployed substructural
language. Its ownership model is essentially global affinity with borrowing.
\Ephapax{} differs in two key respects: (1) per-binding qualifier choice allows
linear bindings that \emph{must} be consumed, whereas Rust's \lstinline{Drop}
trait provides a default destructor; (2) borrowing in \Ephapax{} is explicit
(\lstinline{&x}) rather than implicit through reference creation. The
\lstinline{#[must_use]} attribute in Rust is advisory (a warning, not an error)
and applies to types rather than bindings.

\paragraph{Linear Haskell.}
\citet{bernardy2018linear} introduced linear types to Haskell via multiplicity
annotations on function arrows. This is a different design point from \Ephapax{}:
linearity is a property of how a function \emph{uses} its argument, not a
property of the binding itself. A \lstinline{let}-binding in Linear Haskell is
always unrestricted; linearity flows through the calling convention. \Ephapax{}'s
per-binding linearity is more direct but less compositional.

\paragraph{ATS and Dependent Linear Types.}
\citet{xi2017applied} combined dependent types with linear types in ATS,
achieving safe manual memory management through ``views'' (linear propositions).
ATS separates the proof level (linear) from the value level (unrestricted),
which is a form of dyadic design but at a different granularity---the proof/value
stratification rather than the binding level. \Ephapax{}'s approach is
syntactically simpler but less expressive in the proof dimension.

\paragraph{Quantitative Type Theory.}
QTT~\citep{atkey2018syntax,mcbride2016got} generalizes linearity to arbitrary
resource semirings. Idris 2~\citep{brady2021idris2} and Granule~\citep{orchard2019quantitative}
implement QTT, subsuming linear ($1$) and affine ($\{0,1\}$) as special cases.
The dyadic principle is less general than QTT but has a simpler programmer-facing
model: two binding forms rather than an algebra of grades. We view dyadic design
as a pragmatic restriction of QTT that trades generality for clarity.

\paragraph{Clean and Uniqueness Types.}
Clean~\citep{barendsen1996uniqueness} uses uniqueness types to enable safe
in-place mutation. A unique value is guaranteed to have no aliases, which is
related to linearity but distinct: uniqueness is about reference count, while
linearity is about usage count. Clean's uniqueness types are global (all unique
values follow the same discipline), which differs from \Ephapax{}'s per-binding
choice.

\paragraph{Substructural Type Systems Survey.}
\citet{walker2005substructural} provides a comprehensive survey of substructural
type systems, organizing them by which structural rules they restrict. The dyadic
principle operates at a different level: it is not a novel point in the
structural-rule lattice but a \emph{design methodology} for combining two points
in the lattice within a single language.

\paragraph{Session Types.}
Session types~\citep{honda1998language} enforce communication protocols on
channels, requiring that messages are sent and received in a prescribed order.
Session-typed channels are inherently linear (a channel endpoint must be used
exactly according to its protocol). \Ephapax{}'s linear bindings could encode
session types, providing another instance of the linear/affine dyad: session
channels as linear, message payloads as affine.

\paragraph{Effect Systems.}
Effect systems~\citep{lucassen1988polymorphic} track computational effects
(I/O, state, exceptions) in the type system. The sync/async dyadic pair
(\Cref{sec:generalization}) is related to effect typing: ``sync'' corresponds
to the pure effect and ``async'' to an I/O-like effect. Algebraic effect
handlers~\citep{plotkin2013handling} provide a more general framework, but
the dyadic principle offers a simpler programmer-facing model for the specific
case of two complementary disciplines.

% ============================================================================
% 11. DISCUSSION
% ============================================================================
\section{Discussion}
\label{sec:discussion}

\subsection{Limitations}

The dyadic principle as presented has several limitations:

\begin{enumerate}
  \item \textbf{Binary restriction.} We consider only pairs of disciplines.
    Some design tensions may involve three or more dimensions (e.g., linear,
    affine, and relevant types form a triple). Extending the principle to
    $n$-adic design is straightforward in concept but may complicate the
    programmer model.

  \item \textbf{Qualifier inference.} \Ephapax{} requires explicit qualifier
    annotations (\lstinline{let} vs.\ \lstinline{let!}). A practical
    implementation might benefit from qualifier inference, defaulting to affine
    and inferring linearity where resource obligations are detected. This is
    engineering, not theory, but it affects adoption.

  \item \textbf{Higher-order complications.} When linear and affine closures
    interact---a linear closure capturing an affine binding, or vice versa---the
    typing rules become more complex. We have addressed this through
    qualifier-annotated arrows ($A \xrightarrow{q} B$) but the full treatment
    of higher-order dyadic typing requires further investigation.

  \item \textbf{Concurrency.} The current type system is sequential. Extending
    \Ephapax{} to concurrent programs requires ensuring that linear bindings are
    not shared between threads (since sharing would violate single-use). This is
    analogous to Rust's \lstinline{Send}/\lstinline{Sync} traits and is future
    work.
\end{enumerate}

\subsection{Philosophical Implications}

The dyadic principle suggests a broader lesson about language design: the most
productive design points may not be the ``purest'' ones. A language that is
purely functional, or purely imperative, or purely linear, or purely lazy,
optimizes for one dimension at the expense of others. The dyadic principle
argues for \emph{principled impurity}: combining two disciplines in a way that
is formally sound and mutually beneficial, rather than seeking purity in one
dimension.

This echoes Girard's observation that linear logic is not a replacement for
classical logic but a \emph{refinement} that makes the structural rules
visible~\citep{girard1987linear}. Similarly, the dyadic principle does not
replace monodisciplinary design but refines it, making the tension between
disciplines visible and programmable.

\subsection{The Cost of Choice}

Every binding in \Ephapax{} requires a choice: \lstinline{let} or
\lstinline{let!}. This is a cognitive cost. We argue it is a \emph{beneficial}
cost: it forces the programmer to consider, at the point of introduction,
whether a value carries a resource obligation. This consideration is already
present in the programmer's mind when writing code in any language---the
difference is that \Ephapax{} makes it explicit and checkable.

In our experience, the ratio is approximately 80/20: roughly 80\% of bindings
are affine (data, intermediate computations, formatting results) and 20\% are
linear (file handles, connections, transactions, region handles). The linear
annotations thus serve as \emph{documentation by the programmer, enforced by
the compiler}: they highlight the 20\% of bindings that carry real obligations.

% ============================================================================
% 12. CONCLUSION
% ============================================================================
\section{Conclusion}
\label{sec:conclusion}

We have introduced the \emph{dyadic design principle}, a methodology for
programming language architecture in which two fundamentally different but
inseparable disciplines co-exist at the binding level. We formalized three
constitutive properties---non-diminishment, emergence, and inseparability---that
distinguish dyadic design from both monodisciplinary languages and unprincipled
multi-paradigm languages.

We instantiated the principle in \Ephapax{}, where linear (\lstinline{let!})
and affine (\lstinline{let}) bindings cohabit per-expression. The formal type
system guarantees that linear bindings are consumed exactly once and affine
bindings at most once, with neither discipline interfering with the other. The
interaction of linear region handles with affine interior values yields
region-based memory management without garbage collection---an emergent property
that neither discipline achieves alone.

The session-sentinel case study demonstrated that the dyadic approach prevents
an entire class of resource-leak failures (56 SIGABRT crashes) that
monodisciplinary affine designs admit. The fix is zero-cost: linearity is
enforced at compile time, with no runtime overhead.

Looking forward, we sketched three further dyadic pairs---eager/lazy, sync/async,
and deterministic/probabilistic---arguing that the principle generalizes to any
design tension that is currently resolved by a global architectural commitment.
The dyadic principle offers a different resolution: \emph{don't choose; let the
programmer choose, per binding, and let the type system enforce both choices
simultaneously.}

Like the mother dolphin and her calf, the two disciplines swim together: never
less than one, together more than two.

% ============================================================================
% REFERENCES
% ============================================================================
\bibliographystyle{plainnat}

\begin{thebibliography}{20}

\bibitem[Atkey(2018)]{atkey2018syntax}
Robert Atkey.
\newblock Syntax and semantics of quantitative type theory.
\newblock In \emph{Proceedings of the 33rd Annual ACM/IEEE Symposium on Logic in Computer Science (LICS)}, pages 56--65, 2018.

\bibitem[Barendsen and Smetsers(1996)]{barendsen1996uniqueness}
Erik Barendsen and Sjaak Smetsers.
\newblock Uniqueness typing for functional languages with graph rewriting semantics.
\newblock \emph{Mathematical Structures in Computer Science}, 6(6):579--612, 1996.

\bibitem[Bernardy et~al.(2018)]{bernardy2018linear}
Jean-Philippe Bernardy, Mathieu Boespflug, Ryan~R. Newton, Simon~Peyton Jones, and Arnaud Spiwack.
\newblock Linear {H}askell: practical linearity in a higher-order polymorphic language.
\newblock \emph{Proceedings of the ACM on Programming Languages}, 2(POPL):1--29, 2018.

\bibitem[Brady(2021)]{brady2021idris2}
Edwin Brady.
\newblock \emph{Idris 2: Quantitative Type Theory in Practice}.
\newblock In \emph{European Conference on Object-Oriented Programming (ECOOP)}, 2021.

\bibitem[Girard(1987)]{girard1987linear}
Jean-Yves Girard.
\newblock Linear logic.
\newblock \emph{Theoretical Computer Science}, 50(1):1--102, 1987.

\bibitem[Honda et~al.(1998)]{honda1998language}
Kohei Honda, Vasco~T. Vasconcelos, and Makoto Kubo.
\newblock Language primitives and type discipline for structured communication-based programming.
\newblock In \emph{European Symposium on Programming (ESOP)}, pages 122--138. Springer, 1998.

\bibitem[Jung et~al.(2018)]{jung2018rustbelt}
Ralf Jung, Jacques-Henri Jourdan, Robbert Krebbers, and Derek Dreyer.
\newblock {RustBelt}: Securing the foundations of the {Rust} programming language.
\newblock \emph{Proceedings of the ACM on Programming Languages}, 2(POPL):1--34, 2018.

\bibitem[Lucassen and Gifford(1988)]{lucassen1988polymorphic}
John~M. Lucassen and David~K. Gifford.
\newblock Polymorphic effect systems.
\newblock In \emph{Proceedings of the 15th ACM SIGPLAN-SIGACT Symposium on Principles of Programming Languages (POPL)}, pages 47--57, 1988.

\bibitem[McBride(2016)]{mcbride2016got}
Conor McBride.
\newblock I got plenty o' nuttin'.
\newblock In \emph{A List of Successes That Can Change the World}, pages 207--233. Springer, 2016.

\bibitem[Nystrom(2015)]{nystrom2015colored}
Bob Nystrom.
\newblock What color is your function?
\newblock Blog post, \url{https://journal.stuffwithstuff.com/2015/02/01/what-color-is-your-function/}, 2015.

\bibitem[Orchard et~al.(2019)]{orchard2019quantitative}
Dominic Orchard, Vilem-Benjamin Liepelt, and Harley Eades~III.
\newblock Quantitative program reasoning with graded modal types.
\newblock \emph{Proceedings of the ACM on Programming Languages}, 3(ICFP):1--30, 2019.

\bibitem[Plotkin and Pretnar(2013)]{plotkin2013handling}
Gordon~D. Plotkin and Matija Pretnar.
\newblock Handling algebraic effects.
\newblock \emph{Logical Methods in Computer Science}, 9(4), 2013.

\bibitem[Walker(2005)]{walker2005substructural}
David Walker.
\newblock Substructural type systems.
\newblock In \emph{Advanced Topics in Types and Programming Languages}, chapter~1. MIT Press, 2005.

\bibitem[Weihs(2004)]{weihs2004mechanics}
Daniel Weihs.
\newblock The hydrodynamics of dolphin drafting.
\newblock \emph{Journal of Biology}, 3(2):8, 2004.

\bibitem[Xi(2017)]{xi2017applied}
Hongwei Xi.
\newblock Applied type system: An approach to practical programming with theorem proving.
\newblock \emph{Journal of Functional Programming}, 27:e17, 2017.

\end{thebibliography}

% ============================================================================
% APPENDIX
% ============================================================================
\appendix

\section{Full Typing Rules}
\label{app:typing-rules}

We provide the complete set of typing rules for \Ephapax{}, including
pattern matching, region operations, and recursive functions.

\begin{figure}[htbp]
\centering
\begin{mathpar}
  \inferrule*[right=\textsc{Unit}]
    {~}
    {\judge{\emptyset}{()}{()}}

  \inferrule*[right=\textsc{Int}]
    {n \in \mathbb{Z}}
    {\judge{\emptyset}{n}{\text{Int}}}

  \inferrule*[right=\textsc{Bool}]
    {b \in \{\text{true}, \text{false}\}}
    {\judge{\emptyset}{b}{\text{Bool}}}

  \inferrule*[right=\textsc{Pair}]
    {\judge{\ctx_1}{e_1}{A} \\ \judge{\ctx_2}{e_2}{B}}
    {\judge{\ctx_1 \bowtie \ctx_2}{(e_1, e_2)}{A \times B}}

  \inferrule*[right=\textsc{Pair-Elim}]
    {\judge{\ctx_1}{e}{A \times B} \\
     \judge{\ctx_2, x :^{q_1} A, y :^{q_2} B}{e'}{C}}
    {\judge{\ctx_1 \bowtie \ctx_2}
           {\mathbf{match}\; e\; \mathbf{with}\; (x, y) \Rightarrow e'}{C}}

  \inferrule*[right=\textsc{Region-Create}]
    {~}
    {\judge{\emptyset}{\text{region::create}()}{\text{Region}^\lmark}}

  \inferrule*[right=\textsc{Region-Alloc}]
    {\judge{\ctx_1}{r}{\&\text{Region}} \\
     \judge{\ctx_2}{e}{A}}
    {\judge{\ctx_1 \bowtie \ctx_2}
           {\text{region::alloc}(r, e)}
           {\text{Ref}^\amark\langle A \rangle}}

  \inferrule*[right=\textsc{Region-Destroy}]
    {\judge{\ctx}{r}{\text{Region}^\lmark}}
    {\judge{\ctx}{\text{region::destroy}(r)}{()}}

  \inferrule*[right=\textsc{Fix}]
    {\judge{\ctx, f :^\amark (A \xrightarrow{q} B), x :^q A}{e}{B}}
    {\judge{\ctx}{\mathbf{fix}\; f\; (x :^q A) = e}{A \xrightarrow{q} B}}
\end{mathpar}
\caption{Extended typing rules: primitives, pairs, regions, and recursion.}
\label{fig:extended-rules}
\end{figure}

\section{Metatheory Details}
\label{app:metatheory}

\begin{lemma}[Substitution]
\label{lem:substitution}
If $\judge{\ctx, x :^q A}{e}{B}$ and $\judge{\ctx'}{v}{A}$ and $v$ is a value,
then $\judge{\ctx \bowtie \ctx'}{e[v/x]}{B}$.
\end{lemma}

\begin{proof}
By induction on the derivation of $\judge{\ctx, x :^q A}{e}{B}$.

\textit{Case} \textsc{Var}: $e = x$, $B = A$, $\ctx = \emptyset$.
Then $e[v/x] = v$ and $\judge{\ctx'}{v}{A}$ by assumption.

\textit{Case} \textsc{Var} ($e = y \neq x$): $y :^{q'} B \in \ctx$.
Then $e[v/x] = y$ and $\judge{\ctx}{y}{B}$ with $x :^q A$ removed.
The context $\ctx \bowtie \ctx'$ is well-formed because $x$ no longer appears
and $\ctx'$ is used for the (now unnecessary) substitution.

\textit{Case} \textsc{Let-Affine}: $e = \letaff\; y = e_1\; \mathbf{in}\; e_2$.
By the IH, substitution preserves typing in both sub-derivations.
The context split distributes $x :^q A$ to whichever sub-context contains $x$,
and substitution replaces $x$ with $v$ in the appropriate sub-term.

The remaining cases (\textsc{Let-Linear}, \textsc{Abs}, \textsc{App},
\textsc{Branch}, \textsc{Weaken}, \textsc{Borrow}) follow analogously.
\end{proof}

\begin{lemma}[Context Split Determinacy]
\label{lem:split-det}
For a well-typed term $e$ with context $\ctx$, if $e$ has sub-terms $e_1, e_2$
with $\ctx = \ctx_1 \bowtie \ctx_2$, then the split is uniquely determined by
$\text{FV}(e_1)$ and $\text{FV}(e_2)$ up to weakening of affine bindings.
\end{lemma}

\begin{proof}
Each linear binding $x :^\lmark A$ appears in exactly one of $\text{FV}(e_1)$
or $\text{FV}(e_2)$ (since it must be used exactly once in total), determining
its placement. Each affine binding $x :^\amark A$ appears in at most one of
$\text{FV}(e_1)$ or $\text{FV}(e_2)$, determining its placement (or weakening
if it appears in neither). \qed
\end{proof}

\section{Session-Sentinel Crash Log Excerpts}
\label{app:crash-logs}

For completeness, we provide representative crash signatures from the
session-sentinel failure (anonymized):

\begin{lstlisting}[basicstyle=\ttfamily\scriptsize, frame=single, numbers=none]
session-sentinel[2847]: SIGABRT at dbus_connection_unref+0x42
  Backtrace:
    #0  0x7f3a2b1c4042 in raise (libc.so.6)
    #1  0x7f3a2b1ad4f4 in abort (libc.so.6)
    #2  0x7f3a2c8e1a23 in _dbus_abort (libdbus-1.so.3)
    #3  0x7f3a2c8d2b17 in dbus_connection_unref (libdbus-1.so.3)
    #4  0x5623a41b2f91 in handle_session_event (session-sentinel)
    #5  0x5623a41b1a44 in event_loop (session-sentinel)

session-sentinel[3201]: SIGABRT at dbus_connection_unref+0x42
  Backtrace: [similar pattern, different call site]
  Context: error path in handle_suspend(), connection not closed

session-sentinel[4455]: SIGABRT at dbus_connection_unref+0x42
  Backtrace: [similar pattern]
  Context: timeout handler dropped connection without closing
\end{lstlisting}

All 56 crashes share the same proximate cause:
\lstinline{dbus_connection_unref} is called on a connection that was not first
closed via \lstinline{dbus_connection_close}. The affine semantics of the
original implementation allowed the connection handle to be implicitly dropped
(via scope exit), which triggered \lstinline{unref} without the preceding
\lstinline{close}.

\end{document}
